% Source: https://artofproblemsolving.com/community/c3h1614611_one_sum_on_quadratic_equation
% Stable identity: aops-topic-1614611-q01
Write $A=a^2-5a+3$; for the equation to be quadratic we need $A\ne0$. Let the roots be $r$ and $2r$. By Vieta's formulas,
\[
3r=-\dfrac{3a-1}{A},\qquad 2r^2=\dfrac{2}{A}.
\]
From the second, $r^2=\dfrac{1}{A}$. Squaring the first, $9r^2=\dfrac{(3a-1)^2}{A^2}$, so
\[
\dfrac{9}{A}=\dfrac{(3a-1)^2}{A^2}\quad\Longrightarrow\quad 9A=(3a-1)^2 .
\]
Thus $9a^2-45a+27=9a^2-6a+1$, so $39a=26$ and $a=\dfrac{2}{3}$.

\textbf{Check.} For $a=\dfrac{2}{3}$: $A=\dfrac{4}{9}-\dfrac{10}{3}+3=\dfrac{1}{9}$ and $3a-1=1$, so the equation is $\dfrac{1}{9}x^2+x+2=0$, i.e. $x^2+9x+18=0$, with roots $-3$ and $-6$. One is twice the other. Hence $a=\boxed{\dfrac{2}{3}}$.
